\section{Week 6 -- AI in Finance \& NLP Fundamentals}

\subsection{Foundations of Logic \& Reasoning}
\begin{itemize}
    \item \textbf{Deductive Logic:} Leads down; conclusion is guaranteed to be true.
    \item \textbf{Inductive Logic:} Leads into; conclusion is probably true.
    \item \textbf{Abductive Logic:} Leads away; conclusion is a best guess.
\end{itemize}

\subsection{The 5 Pillars of AI}
\begin{itemize}
    \item \textbf{Supervised:} Task-driven; predicts the next value based on labeled data.
    \item \textbf{Unsupervised:} Data-driven; identifies underlying clusters and patterns.
    \item \textbf{Reinforcement:} Utility-driven; learns from mistakes via trial and error.
    \item \textbf{Transfer:} Reuse-driven; adapts a model trained on Task A for Task B.
    \item \textbf{Deep Learning:} Self-learning; utilizes artificial neural networks.
\end{itemize}

\subsection{Large Language Models (LLMs) Training Pipeline}
\begin{enumerate}
    \item \textbf{Supervised Fine Tuning (SFT):} Collect demonstration data to train a supervised
          policy. Human labelers demonstrate desired output behaviors from prompt
          datasets.
    \item \textbf{Reward Model (RM):} Collect comparison data and train a reward model. Labelers
          rank several model outputs from best to worst.
          \begin{equation*}
              \text{Reward} = f(\text{prompts}, \text{responses}) \text{ }
          \end{equation*}
    \item \textbf{Reinforcement Learning:} Optimize the policy against the RM using Proximal
          Policy Optimization (PPO). The model generates outputs, the RM calculates a
          reward, and the policy updates to maximize rewards.
\end{enumerate}

\subsection{Ethical AI \& Governance}
\begin{itemize}
    \item \textbf{The Dark Side of AI:} Models can produce unfair outcomes, amplify data biases,
          operate as a ``black box,'' and be difficult to explain.
    \item \textbf{Regulatory Frameworks:} Solutions demand equitable decisions (managing impact
          and scope), transparency (explainability and trade-offs), and evolvability
          (balancing risks/rewards and human inputs).
    \item \textbf{Al Automation Framework:} Ranges from ``Human-in-the-loop'' (user makes final
          decision) to ``Human-over-the-loop'' (supervisory role) to
          ``Human-out-of-the-loop'' (AI makes final decision autonomously).
\end{itemize}

\subsection{Natural Language Processing (NLP) Basics}
\begin{itemize}
    \item \textbf{Tokenization:} Transforming text into a series of word tokens, typically
          removing punctuation.
    \item \textbf{POS Tagging:} Attaching parts of speech (nouns, verbs) based on definitions and
          adjacent context.
    \item \textbf{Lemmatization:} Reducing words down to their base or dictionary form (lemma).
    \item \textbf{N-grams:} Adjacent word tokens analyzed together (unigram, bigram, trigram).
    \item \textbf{Sentiment Analysis:} Quantifying subjectivity (opinions, emotional tone) using
          Lexicon-based approaches (manual/dictionary) or Machine Learning.
    \item \textbf{Topic Modelling:} An unsupervised ML method that summarizes vast quantities of
          text to find main ideas. It inputs a document collection and outputs
          distributions of $k$ topics ($Z_1 \dots Z_k$) across each document.
\end{itemize}